Für eine feste Gasmenge gilt die allgemeine Gasgleichung
\begin{align}
 \frac{p_1V_1}{T_1} &= \frac{p_2V_2}{T_2}
\end{align}
Mit $T_1 = \TEiF = \TEi$ folgt
\begin{align}
 T_2 &= \TZwF \\
   &= \TEi\cdot\frac{\pZw\cdot\VZw}{\pEi\cdot\VEi} \\
   &= \TZw
\end{align}
Das entspricht
\begin{align}
 \vartheta_2 &= \thZwF = \thZw \approx \result{\thZwP{0}{3}}
\end{align}
Die Temperatur steigt also um
\begin{align}
 \Delta T &= \dTZF = \dTZ \approx \result{\dTZP{0}{2}}
\end{align}
